\documentclass[11pt]{article}

\usepackage[a4paper,margin=1in]{geometry}
\usepackage[T1]{fontenc}
\usepackage[utf8]{inputenc}
\usepackage{lmodern}
\usepackage{graphicx}
\usepackage{hyperref}

\hypersetup{
  colorlinks=true,
  linkcolor=black,
  urlcolor=blue,
  pdftitle={Telemetry Layer and HAOS Parity Bridge},
  pdfauthor={OpenAI Codex},
  pdfsubject={Stability Oracle telemetry architecture update}
}
\setcounter{secnumdepth}{3}
\setcounter{tocdepth}{2}

\title{Telemetry Layer And HAOS Parity Bridge\\Project Paper}
\author{Stability Oracle Architectural Update}
\date{March 25, 2026}

\begin{document}
\maketitle
\tableofcontents
\newpage

\section{Purpose}

This paper documents the next safe architectural step in Stability Oracle:

\begin{itemize}
\item a parallel telemetry ingestion layer
\item a deterministic HAOS adapter
\item a deterministic temporal normalization surface
\item a geometry encoder v0
\item a narrow bridge back into the current logical engine
\item a golden parity test against the existing logical classification path
\end{itemize}

This is not a rewrite of the current oracle.
It is an additive path that preserves the present logical kernel while preparing a future domain-agnostic telemetry surface.

\section{Architectural Stance}

The project now distinguishes between two layers of concern.

\subsection{Current Stable Core}

The current stable logical kernel remains:

\begin{itemize}
\item \texttt{State}
\item \texttt{Perturbation}
\item metric registry
\item classifier policy
\item router
\item \texttt{OracleEngine}
\end{itemize}

This remains the execution kernel for structural stability evaluation.

\subsection{New Parallel Telemetry Layer}

The new telemetry path sits beside the current core:

\begin{itemize}
\item \texttt{TelemetryFrame}
\item \texttt{TelemetrySequence}
\item HAOS telemetry adapter
\item temporal normalizer
\item geometry encoder v0
\item telemetry bridge
\end{itemize}

This layer is designed to ingest domain-specific records without changing the semantics of the current logical engine.

\begin{figure}[h]
  \centering
  \includegraphics[width=\textwidth]{../../Images/Stability Oracle Evaluation Layer Overview 2.png}
  \caption{Updated public overview for the Stability Oracle evaluation stack.}
\end{figure}

\section{Core Principle}

The telemetry layer follows one rule:

\begin{quote}
The engine does not know what a system means. It only sees the geometry of change.
\end{quote}

This means:

\begin{itemize}
\item no domain-specific semantics in the telemetry core
\item no interpretation of feature labels inside the geometry or metric layers
\item no learned or adaptive behavior in v0
\end{itemize}

\section{Telemetry Contracts}

\subsection{TelemetryFrame}

\texttt{TelemetryFrame} is the minimal domain-agnostic telemetry atom.
It contains:

\begin{itemize}
\item a finite timestamp
\item an entity identifier
\item a finite numeric state vector
\item JSON-friendly metadata
\end{itemize}

\subsection{TelemetrySequence}

\texttt{TelemetrySequence} is the ordered assessment unit.
Its contract enforces:

\begin{itemize}
\item nondecreasing timestamps
\item a consistent entity identifier
\item fixed feature width across all frames in the sequence
\item finite numeric state vectors
\end{itemize}

This distinction matters:
raw upstream sources may have arbitrary dimensionality across domains, but once a sequence is normalized for assessment, the emitted stream must be fixed-width.

\section{HAOS Adapter}

\subsection{Role}

The HAOS telemetry adapter exists to prove that a parallel telemetry path can be built without changing the logical oracle semantics.

\subsection{v0 Output Shape}

For a state pair:

\begin{itemize}
\item frame 0 represents the before-state
\item frame 1 represents the after-state
\end{itemize}

The adapter emits only structural features already allowed by the current logical path, such as:

\begin{itemize}
\item node count
\item edge count
\item feature count
\item timestamp count
\item connected-component count
\item edge density
\item mean degree
\item max degree
\end{itemize}

Optionally, it can append deterministic before/after delta features.

\subsection{Non-Goals}

The adapter does not interpret meaning.
It does not say what a node, edge, timestamp, or feature \emph{means}.
It only converts already-permitted structure into numeric telemetry.

\section{Temporal Normalization}

The temporal normalizer is intentionally minimal in v0.
Its job is to:

\begin{itemize}
\item validate telemetry
\item optionally scale vectors per sequence
\item optionally compute first differences
\item optionally mark simple time gaps in metadata
\end{itemize}

It does not:

\begin{itemize}
\item learn interpolation behavior
\item resample with invented values
\item adapt online
\item introduce stochastic preprocessing
\end{itemize}

Default behavior is close to a no-op apart from deterministic validation and scaling.

\section{Geometry Encoder v0}

The geometry encoder is also intentionally small.
Its permitted operations are:

\begin{itemize}
\item deterministic scaling
\item first derivatives
\item second derivatives if explicitly enabled
\end{itemize}

It does not attempt:

\begin{itemize}
\item learned embeddings
\item adaptive manifold projections
\item ontology-sensitive geometry
\end{itemize}

Its role is only to make change geometry more explicit while remaining deterministic and inspectable.

\section{The Bridge}

The bridge is the key safety constraint.

Instead of forcing the current logical engine to consume arbitrary telemetry directly, the project adds a narrow bridge:

\begin{verbatim}
telemetry_to_state_transition(...)
\end{verbatim}

This bridge only accepts HAOS-shaped telemetry sequences produced by the adapter and only when they contain explicit bridge metadata.
If that metadata is missing or the shape is unsafe, the bridge fails closed.

That decision preserves the existing logical engine as the execution kernel.

\section{Golden Parity Test}

The most important new test is the parity bridge test.

It compares two paths:

\subsection{Direct Logical Path}

\begin{verbatim}
before + after -> OracleEngine.evaluate_transition(...)
\end{verbatim}

\subsection{Telemetry Path}

\begin{verbatim}
before + after
-> HaosTelemetryAdapter
-> TemporalNormalizer
-> StateGeometryEncoder
-> telemetry_to_state_transition(...)
-> OracleEngine.evaluate_transition(...)
\end{verbatim}

Current requirement:

\begin{itemize}
\item classification band parity on canonical fixtures
\item \texttt{stable == stable}
\item \texttt{marginal == marginal}
\item \texttt{unstable == unstable}
\end{itemize}

Exact metric equality is not required at this stage.
Band parity is the safe migration gate.

\section{Why This Matters}

This design prevents a common architectural mistake:
rewriting a working logical oracle into a generic telemetry engine in one jump.

Instead, the project now follows a survivable path:

\begin{enumerate}
\item preserve the current logical core
\item build a parallel telemetry ingestion layer
\item prove parity on HAOS-shaped inputs
\item only then widen the system to non-HAOS domains
\end{enumerate}

\section{First Non-HAOS Demo Direction}

The intended first non-HAOS demo after parity is LLM reasoning trace stability.

This domain is suitable because it is:

\begin{itemize}
\item sequential
\item easy to serialize
\item visible to agents and users
\item free of physics claims
\item useful for monitoring tool-use and chain quality
\end{itemize}

But this paper records a strict sequencing rule:
no non-HAOS generalization should proceed until HAOS parity is established and retained.

\section{What This Update Does Not Do}

This update does not:

\begin{itemize}
\item rewrite the current oracle into a universal time-series engine
\item change metric semantics
\item introduce learned encoders
\item add streaming behavior
\item calibrate recovery as a probability
\item widen domain support beyond a safe parallel contract
\end{itemize}

The term used for downstream scoring remains \texttt{recovery\_score}, not probability language.

\section{Operational Summary}

The repository now supports a deterministic telemetry path such as:

\begin{verbatim}
seq = HaosTelemetryAdapter().from_state_pair(before, after)
seq = TemporalNormalizer().normalize(seq)
seq = StateGeometryEncoder().encode(seq)
before_state, after_state = telemetry_to_state_transition(seq)
\end{verbatim}

and then:

\begin{verbatim}
result = build_default_engine().evaluate_transition(before_state, after_state)
\end{verbatim}

This makes telemetry ingestion parallel to the existing logical path rather than competitive with it.

\section{Conclusion}

The Stability Oracle project now has:

\begin{itemize}
\item a deterministic telemetry contract
\item a deterministic HAOS adapter
\item a minimal temporal normalizer
\item a minimal geometry encoder
\item a narrow telemetry bridge
\item a golden parity test against the existing logical engine
\end{itemize}

This is the right level of change.
It widens the system carefully while preserving the thing that already works.

\end{document}
